\subsection{DEFINITION OF K'-STRUCTURES}

PROPOSITION 1. Let V be a right vector space over K and V' a subset of V which is a vector subspace over K'. The following conditions are equivalent:

\begin{enumerate}
\def\labelenumi{(\alph{enumi})}
\item
  The K-linear mapping \(\lambda\) of \(\mathrm{V}_{(\mathbb{K})}^{\prime} = \mathrm{V}^{\prime}\otimes_{\mathbb{K}^{\prime}}\mathbf{K}\) into V such that \(\lambda (x^{\prime}\otimes \xi) = x^{\prime}\xi\) for all \(x^{\prime}\in \mathrm{V}^{\prime},\xi \in \mathbf{K}\) , is bijective.
\item
  Every K'-linear mapping \(f'\) of V' into a vector K-space W can be extended uniquely to a K-linear mapping f of V into W.
\item
  Every basis of \(\mathbf{V}'\) over \(\mathbf{K}'\) is a basis of \(\mathbf{V}\) over \(\mathbf{K}\) .
\item
  There exists a basis of \(\mathbf{V}'\) over \(\mathbf{K}'\) which is also a basis of \(\mathbf{V}\) over \(\mathbf{K}\) .
\item
  The vector K-space V is generated by V' and every subset of V' free over K' is free over K.
\end{enumerate}

We know (§ 5, no. 1) that \(\mathbf{V}_{(\mathbf{K})}^{\prime}\) has a right vector K-space structure under which \((x^{\prime} \otimes \xi)\eta = x^{\prime} \otimes (\xi\eta)\) ( \(\xi, \eta\) in K, \(x^{\prime} \in \mathbf{V}^{\prime}\) ) and that for every K'-linear mapping \(f^{\prime}\) of V' into a vector K-space W, there exists one and only one K-linear mapping \(\bar{f}^{\prime}\) of \(\mathbf{V}_{(\mathbf{K})}^{\prime}\) into W such that \(\bar{f}^{\prime}(x^{\prime} \otimes 1) = f^{\prime}(x^{\prime})\) for \(x^{\prime} \in \mathbf{V}^{\prime}\) . If \(j\) is the canonical injection of V' into V, \(\lambda\) is just the corresponding K-linear mapping \(\bar{j}\) . If \(\lambda\) is bijective, then for every K'-linear mapping \(f^{\prime}: \mathbf{V}^{\prime} \to \mathbf{W}, \bar{f}^{\prime} \circ \lambda^{-1}\) is the unique K-linear mapping of V into W extending \(f^{\prime}\) ; in other words, (a) implies (b). Conversely, if (b) holds, there exists in particular a K-linear mapping \(\mu\) of V into \(\mathbf{V}_{(\mathbf{K})}^{\prime}\) such that \(\mu(x^{\prime}) = x^{\prime} \otimes 1\) for all \(x^{\prime} \in \mathbf{V}^{\prime}\) ; it is immediate that \(\mu \circ \lambda = 1_{\mathbf{V}_{(\mathbf{K})}^{\prime}}\) ; on the other hand, \(\lambda(\mu(x^{\prime})) = x^{\prime}\) for all \(x^{\prime} \in \mathbf{V}^{\prime}\) and, as by hypothesis \(j: \mathbf{V}^{\prime} \to \mathbf{V}\) can be extended uniquely to an endomorphism of V, of necessity \(\lambda \circ \mu = 1_{\mathbf{V}}\) , which completes the proof that (a) and (b) are equivalent.

For every basis \(\mathbf{B}'\) of \(\mathbf{V}'\) over \(\mathbf{K}'\) , the set \(\mathbf{B}\) of elements \(v' \otimes 1\) of \(\mathbf{V}_{(\mathbf{K})}'\) , where \(v'\) runs through \(\mathbf{B}'\) , is a basis of \(\mathbf{V}_{(\mathbf{K})}'\) over \(\mathbf{K}\) ( \(\S 5\) , no. 1, Proposition 4) and \(\lambda(\mathbf{B}) = \mathbf{B}'\) . For \(\lambda\) to be bijective, it is necessary that the image under \(\lambda\) of every basis of \(\mathbf{V}_{(\mathbf{K})}'\) be a basis of \(\mathbf{V}\) over \(\mathbf{K}\) and it is sufficient that it be so for a single basis of \(\mathbf{V}_{(\mathbf{K})}'\) ( \(\S 1\) , no. 11, Corollary 2 to Proposition 17). This proves the equivalence of (a), (c) and (d).

As every subset of \(\mathbf{V}'\) which is free over \(\mathbf{K}'\) is contained in a basis of \(\mathbf{V}'\) over \(\mathbf{K}'\) (§ 7, no. 1, Theorem 2), (c) implies (e). Finally, suppose that (e) holds; if \(\mathbf{B}'\) is a basis of \(\mathbf{V}'\) over \(\mathbf{K}'\) , it is a free subset of \(\mathbf{V}\) over \(\mathbf{K}\) ; on the other hand \(\mathbf{B}'\) generates \(\mathbf{V}'\) over \(\mathbf{K}'\) and hence generates \(\mathbf{V}\) over \(\mathbf{K}\) by hypothesis; therefore \(\mathbf{B}'\) is a basis of \(\mathbf{V}\) over \(\mathbf{K}\) , which proves that (e) implies (c).

DEFINITION 1. Let V be a right vector space over a field K and K' a subfield of K. Every vector sub-K'-space V' of V satisfying the equivalent conditions of Proposition 1 is called a K'-structure on V.

Example. Let B be a basis of V over K. For every subfield \(K'\) of K, the vector sub- \(K'\) -space of V generated by B admits B as a basis over \(K'\) and hence is a \(K'\) -structure on V. *For example, if K is commutative and V is taken to be the polynomial K-algebra \(K[X_{1},\ldots,X_{n}]\) , then, for every subfield \(K'\) of K, \(K'[X_{1},\ldots,X_{n}]\) is a \(K'\) -structure on \(V\).*

